\chapter{Food Allergies in Children}

\section*{Definition and Overview}
Food allergy in children is an immune reaction to a food, most common in the first years of life. Australian infants have among the highest rates of food allergy in the world, with egg, peanut, cow's milk, and wheat the most common triggers. Reactions range from mild skin and gut symptoms to severe anaphylaxis. Many childhood food allergies, particularly to egg and milk, are outgrown, while others, especially peanut and tree nut allergy, often persist. Diagnosis by a specialist, clear management plans, and careful avoidance keep children safe, and early introduction of common allergens now reduces the risk of developing allergy in the first place. This chapter explains food allergy in children, its recognition, and its management.

\section*{Epidemiology}
Around 1 in 10 Australian infants has a food allergy, one of the highest rates worldwide. Egg allergy affects around 9% of infants, peanut allergy around 3%, and cow's milk allergy a similar proportion. Most food allergies appear in the first two years of life, and hospital admissions for food-related anaphylaxis in children have risen steeply in recent decades. Around half of children outgrow egg allergy by school age, and most outgrow milk allergy, while peanut and tree nut allergies persist in the majority. The rising rates have driven research that changed prevention advice: early introduction of allergenic foods now reduces allergy risk.

\section*{Aetiology and Pathophysiology}
\begin{itemize}
    \item \textbf{The immune system:} the child's immune system produces IgE antibodies against a food protein, treating it as a threat
    \item \textbf{Immediate reactions:} on eating the food, histamine and other chemicals are released, causing symptoms within minutes to 2 hours
    \item \textbf{Anaphylaxis:} a severe, generalised reaction with airway swelling, breathing difficulty, and collapse, requiring immediate adrenaline
    \item \textbf{Non-IgE reactions:} delayed gut symptoms, including vomiting, diarrhoea, and reflux, from other immune mechanisms, most often with milk and soy
    \item \textbf{Risk factors:} eczema in early life is the strongest predictor; family history and delayed introduction of allergens also contribute
    \item \textbf{The "allergic march":} children with food allergy often go on to develop asthma, allergic rhinitis, and other allergies
    \item \textbf{Sensitisation versus allergy:} a positive test shows sensitisation; clinical allergy is confirmed by symptoms after eating the food
\end{itemize}

\section*{Clinical Features}
\begin{itemize}
    \item Immediate reactions: hives, swelling of the lips and face, vomiting, abdominal pain, and wheezing
    \item Anaphylaxis: difficulty breathing, throat tightness, hoarse voice, dizziness, and collapse
    \item Skin changes: flushing, itching, and worsening eczema after the food
    \item Non-IgE reactions: vomiting, diarrhoea, colic, poor feeding, and blood in the stool, hours after feeding
    \item The reaction usually occurs within minutes of the first few feeds or exposures
    \item Severity varies: a child can have a mild reaction one time and a severe reaction another
    \item Growth and nutrition can be affected in unrecognised or poorly managed allergy
\end{itemize}

\section*{Differential Diagnosis}
Symptoms after eating have several causes.
\begin{itemize}
    \item Food intolerance, such as lactose intolerance, without immune involvement
    \item Coeliac disease, an autoimmune reaction to gluten
    \item Gastroenteritis and food poisoning
    \item Reflux, colic, and other causes of unsettled feeding in infants
    \item Eczema, which can be confused with or coexist with allergy
    \item A reaction to a contaminant or additive rather than the food itself
    \item Specialist assessment distinguishes these reliably
\end{itemize}

\section*{When to Seek Medical Care}
Parents should seek medical care after any suspected food reaction, and should have the child assessed by a doctor with allergy expertise before reintroducing the food. Children with eczema and a family history of allergy benefit from early allergy review.

\textbf{Call 000 (or local emergency services) for:}
\begin{itemize}
    \item Difficulty breathing, wheezing, or throat tightness after eating
    \item Swelling of the tongue or face, a hoarse voice, or difficulty swallowing
    \item Dizziness, collapse, or loss of consciousness
    \item Repeated vomiting with breathing difficulty
    \item Any suspected anaphylaxis; give the adrenaline autoinjector immediately if available
\end{itemize}

\section*{Diagnosis and Investigations}
\begin{itemize}
    \item \textbf{History:} the food, the amount, the timing, and the symptoms, documented carefully
    \item \textbf{Skin prick testing:} performed by an allergist or trained doctor
    \item \textbf{Blood tests:} specific IgE levels to the suspected foods
    \item \textbf{Oral food challenge:} supervised eating of the food in hospital, the definitive test
    \item \textbf{ASCIA action plan:} developed for confirmed allergy, with prescriptions for adrenaline where needed
    \item \textbf{Reassessment over time:} children with egg and milk allergy are retested, as many outgrow them
\end{itemize}

\section*{Management}
\begin{itemize}
    \item \textbf{Strict avoidance:} the cornerstone, with label reading, safe food preparation, and avoidance of cross-contamination
    \item \textbf{Adrenaline autoinjectors:} for children at risk of anaphylaxis, with training for parents, schools, and carers
    \item \textbf{ASCIA action plan:} a written, individualised plan for recognising and treating reactions, shared with childcare and school
    \item \textbf{Antihistamines:} for mild symptoms only, not for anaphylaxis
    \item \textbf{Dietitian input:} to ensure the elimination diet is nutritionally complete and to guide reintroduction
    \item \textbf{Oral immunotherapy:} available in specialist centres for some children with peanut allergy, building tolerance under supervision
    \item \textbf{Education:} teaching the child from an early age to ask about food and to carry emergency medicines
    \item \textbf{Reintroduction:} supervised reintroduction of foods that may have been outgrown
\end{itemize}

\section*{Complications}
\begin{itemize}
    \item Anaphylaxis, which can be fatal without prompt adrenaline
    \item Poor growth and nutritional deficiencies from unmanaged elimination diets
    \item Anxiety and reduced quality of life for the child and family
    \item Social exclusion at parties, school camps, and events
    \item Bullying related to food allergy
    \item Accidental exposures despite precautions
\end{itemize}

\section*{Prognosis and Outcomes}
The outlook for children with food allergy is good with proper management. Many children outgrow egg and milk allergy by school age, and a proportion outgrow peanut and other allergies; those who do not are managed safely with avoidance and preparedness. Children who are well educated about their allergy, carry adrenaline when needed, and have supportive schools grow into young people who manage their allergy confidently. Oral immunotherapy is expanding options, and the prognosis continues to improve with research and better care.

\section*{Prevention and Self-Care}
\begin{itemize}
    \item Introduce peanut and egg from around 6 months, once solids start, following the national guidelines, which reduce allergy risk
    \item Seek advice before introducing allergens if the child has severe eczema or a family history of allergy
    \item Breastfeed for as long as possible; this supports health, though it does not prevent allergy on its own
    \item Read labels and teach carers, grandparents, and schools about the allergy
    \item Carry the action plan and adrenaline autoinjector everywhere, and check expiry dates
    \item Practise the autoinjector with a training device
    \item Attend allergy reviews and follow-up testing
    \item Seek psychological support for the family if the allergy is causing anxiety
\end{itemize}

\section*{Sources and Further Reading}
The following reputable sources provide further information for healthcare students and parents seeking reliable, current advice about food allergies in children.
\begin{itemize}
    \item Allergy \& Anaphylaxis Australia. \textit{Food allergy in children.} https://allergyfacts.org.au/
    \item Australasian Society of Clinical Immunology and Allergy (ASCIA). \textit{Pacdiatric food allergy.} https://www.allergy.org.au/
    \item National Allergy Council. \textit{Nip allergies in the Bub program.} https://nationalallergycouncil.org.au/
    \item Raising Children Network. \textit{Food allergy in children.} https://raisingchildren.net.au/
    \item NHS. \textit{Food allergies in children.} https://www.nhs.uk/conditions/food-allergy/
    \item World Health Organization. \textit{Infant feeding and allergy.} https://www.who.int/
\end{itemize}
